\documentclass[12pt,onecolumn]{IEEEtran}


\usepackage[utf8]{inputenc}
\usepackage[T1]{fontenc}
\usepackage[spanish,es-tabla]{babel}
\usepackage{amsmath}
\usepackage{graphicx}
\usepackage{hyperref}
\usepackage{cite}
\usepackage{geometry}
\usepackage{float}
\usepackage{booktabs}
\usepackage{array}
\usepackage{siunitx}

\geometry{margin=1in}

\title{Proyecto Final: Semáforo Vehicular y Peatonal}

\author{
Gabriel Alonso Chavarría Rodríguez\\
José Andrés Guerrero Araya\\
Julián Murillo Wo Ching\\
Jimena Vargas Vega
}

\date{\today}

\begin{document}

% The paper headers
\markboth{EL3215 Laboratorio de Electrónica Analógica, IS2026}%
{EL3215 Laboratorio de Electrónica Analógica}

\maketitle

\begin{abstract}
En este trabajo se presenta el diseño, simulación e implementación física de un sistema de semáforo vehicular y peatonal controlado completamente mediante lógica digital discreta, sin el uso de microcontroladores ni FPGAs. El sistema integra seis subsistemas principales: alimentación rectificada regulada a 5~V, generación de reloj de 1~Hz con temporizador NE555, botón peatonal con antirrebote RC y Schmitt Trigger, contador regresivo BCD de 25 segundos con detección combinacional de los instantes de cambio de estado, máquina de estados finitos (FSM) de cuatro estados implementada con flip-flops 74HC74 y compuertas lógicas, y un subsistema de audio con tres osciladores NE555, selección de ritmo por switch analógico TC4066 y amplificación en dos etapas (2N2222 + TIP31C) para excitar un parlante de $8\,\Omega$. Todos los subsistemas fueron verificados en simulación con Multisim y posteriormente construidos e integrados sobre protoboard. Los resultados obtenidos confirman el correcto funcionamiento de la secuencia de estados, la temporización, la señalización auditiva con cambio automático de ritmo, y la visualización del conteo regresivo en displays de siete segmentos.
\end{abstract}

\begin{IEEEkeywords}
Semáforo, FSM, temporización, lógica digital, electrónica analógica.
\end{IEEEkeywords}

%%%%%%%%%%%%%%%%%%%%%%%%%%%%%%%%%%%%%%%%%%%%%%%%%%%%%%%%%%%%%%
\section{Introducción}

Los sistemas de semáforos son una de las aplicaciones más visibles de la electrónica digital en la infraestructura urbana. Más allá de su función aparente, integran múltiples subsistemas que deben coordinarse con precisión: temporización, secuenciación de estados, detección de eventos externos y señalización tanto visual como auditiva. Su diseño representa un caso de estudio completo en sistemas digitales, ya que exige soluciones a problemas reales como el rebote mecánico de pulsadores, la generación de señales de reloj estables y la amplificación de audio para condiciones de uso en exteriores.

El presente proyecto corresponde al diseño e implementación de un semáforo vehicular con cruce peatonal desarrollado como proyecto final del curso EL3215 Laboratorio de Electrónica Analógica (IS2026, Grupo~8). La restricción principal del proyecto establece que el sistema debe implementarse exclusivamente con componentes de lógica digital discreta y electrónica analógica, sin el uso de microcontroladores, FPGAs ni ningún dispositivo programable. Esta restricción obliga a resolver cada función del sistema a nivel de compuertas, flip-flops y componentes analógicos básicos, lo que constituye en sí mismo el objetivo pedagógico central del proyecto.

El sistema diseñado opera con cuatro estados definidos por una máquina de estados finitos: estado inicial con semáforo vehicular en verde ($S_0$), transición con luz amarilla vehicular ($S_1$), paso peatonal activo ($S_2$) y advertencia de cierre de paso con ritmo de audio acelerado ($S_3$). La transición entre estados es controlada por un contador regresivo BCD de 25~s, cuya lógica combinacional genera pulsos en los instantes precisos de cambio. La solicitud de cruce peatonal se introduce mediante un botón con circuito antirrebote y queda memorizada en un flip-flop hasta que la FSM la atiende.

El presente documento describe el diseño, análisis teórico, simulación e implementación en protoboard de cada subsistema, incluyendo los valores finales de componentes utilizados en el prototipo físico.

%%%%%%%%%%%%%%%%%%%%%%%%%%%%%%%%%%%%%%%%%%%%%%%%%%%%%%%%%%%%%%

\section{Subsistema de Alimentación}

El voltaje del sistema será proporcionado en su totalidad por una fuente de \qty{120}{VAC}, de manera que el semáforo funcione al ser conectado a un tomacorrientes común. Esta señal pasará por un transformador de \qty{120}{VAC} a \qty{12}{VAC}. La señal de este transformador luego será procesada por un arreglo rectificador de 4 diodos 1N4007 con un capacitor de filtro de \qty{4700}{\mu F}. El subsistema hasta este punto, produce una señal promedio de \qty{15}{V}. Para obtener la señal final en DC, se utilizaron dos reguladores de voltaje LM7805 de \qty{5}{V} en paralelo. Uno de estos reguladores es responsable de la alimentación del subsistema de audio, mientras que el otro actua como fuente para los subsistemas restantes.

\begin{figure}[H]
    \centering
    \begin{center}\textit{(nueva foto del diagrama)}\end{center}
    \caption{Subsistema de alimentación rectificador}
    \label{fig:bloque1_555}
\end{figure}

Además, para ambos reguladores se conectó un capacitor de \qty{1}{\mu F} entre su entrada y tierra. Esto se llevó a cabo para estabilizar el voltaje de entrada que ve el regulador ante ruido que puede aparecer en el circuito rectificador. Además, protege al regulador de caídas en el voltaje de entrada debido a cambios abruptos en la corriente utilizada por el sistema. Asimismo, un capacitor de \qty{0.1}{\mu F} entre su salida y tierra por las mismas razones. \cite{7805datasheet}

La decisión de utilizar dos reguladores en paralelo emergió tras observar muy altas temperaturas en el regulador cuando solo se utilizó uno para la alimentación del sistema completo. El 7805 tiene una corriente máxima de \qty{1.5}{A} \cite{7805datasheet}. Después de una serie de mediciones, se vió que el subsistema de audio requería alrededor de \qty{800}{mA}, mientras que el resto de los subsistemas juntos requeren alrededor de \qty{600}{mA}. Por motivos de seguridad y preservación del sistema, se decidió entonces utilizar un regulador exclusivamente para la alimentación del subsistema de audio, cuya señal fue denotada $V_{CC}$. La señal del regulador responsable por la alimentación del resto de subsistemas se denotó $V_S$.

\section{Subsistema de Temporización}

Para el funcionamiento adecuado del sistema de control, y en general del proyecto, se requiere una señal de reloj que permita establecer una referencia de tiempo para los contadores y la máquina de estados. Por esta razón, se diseñó un circuito basado en el temporizador NE555 en configuración astable, ya que esta configuración permite generar una señal cuadrada periódica sin necesidad de una señal externa de disparo.

La frecuencia de oscilación del 555 en modo astable se aproxima mediante \cite{electronicafp}:

\begin{equation}
	\label{eq:NE555}
    f = \frac{1.44}{(R_A + 2R_B)C}
\end{equation}

El objetivo fue obtener una frecuencia cercana a $1~Hz$, de manera que cada ciclo de la señal represente aproximadamente un segundo dentro del sistema. Inicialmente se seleccionaron valores comerciales para las resistencias y el capacitor, y posteriormente se ajustó el valor de $R_A$ para acercar la frecuencia al valor deseado.

Los valores utilizados fueron:

\begin{itemize}
    \item $R_A = 12~k\Omega$
    \item $R_B = 68~k\Omega$
    \item $C = 10~\mu F$
\end{itemize}

Sustituyendo estos valores en la ecuación:

\begin{equation}
    f = \frac{1.44}{(12k\Omega + 2(68k\Omega))(10\mu F)}
\end{equation}

\begin{equation}
    f \approx 0.973~Hz
\end{equation}

Por lo tanto, el período aproximado es:

\begin{equation}
    T = \frac{1}{f} \approx 1.03~s
\end{equation}

Este valor se considera adecuado para la propuesta de diseño, ya que se encuentra cercano a $1~Hz$ y permite utilizar la salida del 555 como reloj base del sistema. La salida del temporizador se conecta a un LED únicamente como indicador visual durante la simulación (Figura 1).

\begin{figure}[H]
    \centering
    \begin{center}\textit{(nueva foto del diagrama)}\end{center}
    \caption{Circuito del reloj de aproximadamente 1 Hz implementado con temporizador 555 en configuración astable.}
    \label{fig:bloque1_555}
\end{figure}

\section{Subsistema de Despliegue}

Este subsistema opera en base a dos contadores CMOS CD4029 conectados en cascodo, uno encargado de contar decenas, y el otro unidades. Ambos contadores se encuentran en modo de contado por década, de modo regresivo, y conectados a la señal de reloj del subsistema de temporización. El contador de decenas se tiene preestablecido empezar en 2, y el de unidades en 5, de modo que se inicie una cuenta regresiva a partir de 25 cuando se habilita el pin de \textit{Preset Enable} de ambos contadores.

Las salidas del contador de las decenas, del bit menos significativo al bit más significativo, se definen $D_0$, $D_1$, $D_2$, y $D_3$. Las salidas del contador de unidades se definen $U_0$, $U_1$, $U_2$, y $U_3$, de la misma manera. Ambos contadores tienen un decodificador CD4511B respectivo, al cuál las salidas del contador están conectadas. Ambos decodificadores están conectados a un display de 7 segmentos de cátodo común correspondiente. El cátodo de estos está conectado a tierra mediante una resistencia de \qty{330}{\Omega} para limitar la corriente. 


\section{Subsistema de Audio}

El subsistema de audio tiene como objetivo generar señalización auditiva para el semáforo peatonal. Según las especificaciones del proyecto, el sistema debe producir un sonido periódico mientras el paso peatonal está habilitado, y aumentar la frecuencia de ese sonido durante los últimos cinco segundos antes del cambio al estado de detención.

Para lograr esto se diseñó un subsistema compuesto por tres etapas principales: generación de tono y ritmo mediante temporizadores NE555, selección lógica del ritmo adecuado según la fase del semáforo, y amplificación de la señal para excitar un parlante de $8\,\Omega$.

\subsection{Generación de Tono y Ritmo}

Se utilizaron tres temporizadores NE555 configurados en modo astable~\cite{ne555art}. Esta es la misma configuración que la del 555 en el subsistema de temporización, y permite generar una onda cuadrada de forma autónoma cuya frecuencia depende de los valores de $R_a$, $R_b$ y el capacitor de temporización $C$. Los tres temporizadores cumplen funciones distintas: uno genera el tono audible, uno genera el ritmo lento para el paso peatonal normal, y uno genera el ritmo rápido para los últimos cinco segundos.
\\
\\
\textbf{NE555 de Tono ($\approx 300\,\text{Hz}$).}
Este temporizador genera la frecuencia audible que percibe el peatón. Su salida por el pin~3 produce una onda cuadrada de aproximadamente $300\,\text{Hz}$, representativa de los tonos utilizados en semáforos peatonales reales. Los valores de componentes seleccionados se muestran en la Tabla~\ref{tab:555tono}.

\begin{table}[H]
\centering
\caption{Componentes del NE555 de tono.}
\label{tab:555tono}
\begin{tabular}{lll}
\toprule
\textbf{Componente} & \textbf{Valor} & \textbf{Justificación} \\
\midrule
$R_a$       & $1\,\text{k}\Omega$   & Protege el pin de descarga \\
$R_b$       & $20\,\text{k}\Omega$  & Define la frecuencia junto con $C$ \\
$C$ timing  & $100\,\text{nF}$      & Capacitor de temporización \\
\bottomrule
\end{tabular}
\end{table}

Aplicando la ecuación~\eqref{eq:NE555}:

\begin{equation}
    f = \frac{1{,}44}{(1\,\text{k} + 2\cdot 20\,\text{k})\cdot 100\,\text{nF}}
      = \frac{1{,}44}{41\,000 \times 100\times10^{-9}}
      \approx 351\,\text{Hz}
\end{equation}

El pin~4 (RST) de este temporizador recibe la señal del bloque de selección de ritmo (TC4066BP), lo que permite interrumpir el tono rítmicamente para crear los bips audibles.

\begin{figure}[H]
    \centering
    \begin{center}\textit{(nueva foto del diagrama)}\end{center}
    \caption{Circuito del NE555 de tono ($\approx 351\,\text{Hz}$).}
    \label{fig:tono}
\end{figure}

\textbf{NE555 de Ritmo Lento ($\approx 5{,}7\,\text{Hz}$).}
Este temporizador genera la cadencia de bips durante la fase normal del paso peatonal. Para lograr una pausa notablemente más larga que el bip, se incorporó un diodo 1N4149 en paralelo con $R_b$, separando los tiempos de carga y descarga del capacitor:

\begin{align}
    t_{\text{HIGH}} &= 0{,}693 \cdot R_a \cdot C \\
    t_{\text{LOW}}  &= 0{,}693 \cdot R_b \cdot C
\end{align}

Los valores seleccionados se muestran en la Tabla~\ref{tab:555lento}.

\begin{table}[H]
\centering
\caption{Componentes del NE555 de ritmo lento.}
\label{tab:555lento}
\begin{tabular}{lll}
\toprule
\textbf{Componente} & \textbf{Valor} & \textbf{Función} \\
\midrule
$R_a$  & $4{,}7\,\text{k}\Omega$ & Duración del bip ($\approx 72\,\text{ms}$) \\
$R_b$  & $6{,}8\,\text{k}\Omega$ & Duración de la pausa ($\approx 104\,\text{ms}$) \\
$C$    & $22\,\mu\text{F}$       & Capacitor de temporización \\
Diodo  & 1N4149                & Separa tiempos de carga y descarga \\
\bottomrule
\end{tabular}
\end{table}

El pin~4 (RST) recibe la señal $AUD$, proveniente de la máquina de estados. Mientras esa señal está en HIGH, el temporizador oscila y genera el ritmo; cuando baja a LOW, el temporizador se detiene y el audio se apaga.

\begin{figure}[H]
    \centering
    \begin{center}\textit{(nueva foto del diagrama)}\end{center}
    \caption{Circuito del NE555 de ritmo lento ($\approx 5{,}7\,\text{Hz}$).}
    \label{fig:ritmo_lento}
\end{figure}

\textbf{NE555 de Ritmo Rápido ($\approx 14{,}6\,\text{Hz}$).}
Este temporizador genera la cadencia acelerada durante los últimos cinco segundos del paso peatonal, comunicando urgencia al peatón. Opera con los mismos principios que el de ritmo lento pero con $C$ reducido para aumentar la frecuencia. Los valores se muestran en la Tabla~\ref{tab:555rapido}.

\begin{table}[H]
\centering
\caption{Componentes del NE555 de ritmo rápido.}
\label{tab:555rapido}
\begin{tabular}{lll}
\toprule
\textbf{Componente} & \textbf{Valor} & \textbf{Función} \\
\midrule
$R_a$  & $1\,\text{k}\Omega$    & Duración del bip ($\approx 3\,\text{ms}$) \\
$R_b$  & $20\,\text{k}\Omega$  & Duración de la pausa ($\approx 65\,\text{ms}$) \\
$C$    & $4{,}7\,\mu\text{F}$  & Capacitor de temporización \\
Diodo  & 1N4149                & Separa tiempos de carga y descarga \\
\bottomrule
\end{tabular}
\end{table}

\begin{figure}[H]
    \centering
    \begin{center}\textit{(nueva foto del diagrama)}\end{center}
    \caption{Circuito del NE555 de ritmo rápido ($\approx 14{,}6\,\text{Hz}$).}
    \label{fig:ritmo_rapido}
\end{figure}

\subsection{Selección Lógica del Ritmo}

Para que el subsistema de audio cambie automáticamente entre ritmo lento y ritmo rápido según la fase del semáforo, se implementó un bloque de selección lógica compuesto por un switch analógico TC4066BP y un inversor 74LS14N.

Se utilizó uno de los seis inversores Schmitt del 74LS14N para invertir la señal de \textit{últimos 5~s} proveniente de la máquina de estados. Esta señal invertida controla el Switch~1 del TC4066BP, garantizando que ambos switches operen de forma complementaria y nunca estén activos simultáneamente:

\begin{itemize}
    \item Señal \textit{últimos 5~s} en LOW $\Rightarrow$ inversor da HIGH $\Rightarrow$ Switch~1 activo $\Rightarrow$ ritmo lento pasa al tono.
    \item Señal \textit{últimos 5~s} en HIGH $\Rightarrow$ inversor da LOW $\Rightarrow$ Switch~1 inactivo $\Rightarrow$ ritmo rápido toma el control.
\end{itemize}

Se utilizaron dos switches del TC4066BP para seleccionar cuál señal de ritmo llega al pin~4 (RST) del NE555 de tono: Switch~1 con la salida del 555 de ritmo lento, y Switch~2 con la salida del 555 de ritmo rápido controlado directamente por la señal \textit{últimos 5~s}. De esta forma el TC4066BP actúa como un multiplexor de dos entradas cuya señal de control proviene directamente de la máquina de estados.

\subsection{Etapa de Amplificación}

La señal del pin~3 del NE555 de tono tiene una amplitud de aproximadamente $5\,\text{V}$ con $V_{CC}$~=~\qty{5}{V}, suficiente en voltaje pero con poca capacidad de corriente para excitar directamente un parlante de $8\,\Omega$. Por ello se diseñó un amplificador de dos etapas: emisor común para amplificar el voltaje, seguido de un seguidor de emisor para amplificar la corriente.

Antes de ingresar al amplificador, la señal del NE555 se atenúa mediante un divisor de tensión ($47\,\text{k}\Omega$ en serie, con capacitor de acople de $4{,}7\,\mu\text{F}$) para reducirla de $5\,\text{V}$ a aproximadamente $800$--$900\,\text{mV}$, evitando la saturación del transistor de la primera etapa.

\textbf{Primera etapa: Emisor Común (2N2222).}
La polarización se realiza con un divisor de tensión en la base para establecer el punto Q centrado en $V_{CE} \approx V_{CC}/2$ \cite{2n2222ds}. Los componentes se muestran en la Tabla~\ref{tab:emisor_comun}.

\begin{table}[H]
\centering
\caption{Componentes de la primera etapa (emisor común).}
\label{tab:emisor_comun}
\begin{tabular}{lcc}
\toprule
\textbf{Componente} & \textbf{Valor} & \textbf{Justificación} \\
\midrule
$R_1$ (divisor base, $V_{CC}$) & $47\,\text{k}\Omega$    & Fija la corriente del divisor \\
$R_2$ (divisor base, GND)      & $8{,}2\,\text{k}\Omega$ & Establece $V_B \approx 0{,}74\,\text{V}$ \\
$R_C$ (colector)               & $1{,}2\,\text{k}\Omega$ & Define $I_C$ y el punto Q \\
$C$ acople entre etapas        & $100\,\mu\text{F}$      & Bloquea el DC entre etapas \\
\bottomrule
\end{tabular}
\end{table}

El punto de operación Q se calcula como:

\begin{align}
    V_B &= V_{CC} \cdot \frac{R_2}{R_1 + R_2}
         = 5 \cdot \frac{8{,}2\,\text{k}}{55{,}2\,\text{k}}
         \approx 0{,}74\,\text{V} \\
    I_C &= \frac{V_{CC} - V_{CE}}{R_C}
         = \frac{5 - 2{,}5}{1200}
         \approx 2{,}08\,\text{mA} \\
    V_C &\approx 2{,}9\,\text{V} \quad \text{(medido en simulación)}
\end{align}

\textbf{Segunda etapa: Seguidor de Emisor (TIP31C).}
La configuración de colector común con TIP31C amplifica la corriente para excitar el parlante, presentando alta ganancia de corriente ($h_{FE} \geq 25$ con $I_C = 1\,\text{A}$ según datasheet \cite{tip31ds}) y baja impedancia de salida. Los componentes se listan en la Tabla~\ref{tab:seguidor_emisor}.

\begin{table}[H]
\centering
\caption{Componentes de la segunda etapa (seguidor de emisor).}
\label{tab:seguidor_emisor}
\begin{tabular}{lcc}
\toprule
\textbf{Componente} & \textbf{Valor} & \textbf{Justificación} \\
\midrule
$R_6$ (divisor base, $V_{CC}$) & $820\,\Omega$           & Fija la corriente del divisor \\
$R_7$ (divisor base, GND)      & $2{,}2\,\text{k}\Omega$ & Establece $V_B \approx 2{,}5\,\text{V}$ \\
Parlante (carga)               & $8\,\Omega$             & Carga en el emisor del TIP31C \\
\bottomrule
\end{tabular}
\end{table}

Con $V_B = 2{,}5\,\text{V}$, la corriente y potencia entregadas al parlante resultan:

\begin{align}
    V_E &= V_B - 0{,}7\,\text{V} = 1{,}8\,\text{V} \\
    I_C &= \frac{V_E}{R_L} = \frac{1{,}8}{8} = 225\,\text{mA} \\
    P   &= I_C^2 \cdot R_L = (0{,}225)^2 \times 8 \approx 405\,\text{mW}
\end{align}

Este valor de potencia es suficiente para garantizar la audibilidad del semáforo peatonal en condiciones normales de uso.

\begin{figure}[H]
    \centering
    \begin{center}\textit{(nueva foto del diagrama)}\end{center}
    \caption{Esquemático del amplificador de dos etapas (emisor común + seguidor de emisor).}
    \label{fig:amplificador}
\end{figure}

\begin{figure}[H]
    \centering
    \begin{center}\textit{(nueva foto del diagrama)}\end{center}
    \caption{Esquemático completo del subsistema de audio en Multisim.}
    \label{fig:subsistema_audio}
\end{figure}

%%%%%%%%%%%%%%%%%%%%%%%%%%%%%%%%%%%%%%%%%%%%%%%%%%%%%%%%%%%%%%%%%%%%%%%%%%%%%%%%%%%%%%%%%%%%%%%%%%%%%%%%%%%%%%%%%%%%%%%%%%%%%%%%%%%%%%%%%%%%%%%%%%%%%%%%%%%%%

%% PARTE DE LA FSM

%%%%%%%%%%%%%%%%%%%%%%%%%%%%%%%%%%%%%%%%%%%%%%%%%%%%%%%%%%%%%%%%%%%%%%%%%%%%%%%%%%%%%%%%%%%%%%%%%%%%%%%%%%%%%%%%%%%%%%%%%%%%%%%%%%%%%%%%%%%%%%%%%%%%%%%%%%%%%%%%%%%%%%%%%%%%%%%



\section{Máquina de estados (FSM)}



\subsection{Contador y Drivers}

La máquina de estados está basada en un contador CMOS CD4029 sincronizado al reloj del sistema. Este se utilizó por su capacidad de poder moverse entre diferentes estados de manera secuencial y reinicio asincrónico. Se utilizaron sus dos primeras salidas $Q_0$ y $Q_1$, correspondientes a los dos bits menos significativos de la cuenta, para representar los cuatro estados necesarios. En conjunto a una compuera AND CD4081 y una NOT CD4069 se pasó de esta cuenta en binario de dos bits a una condificación \textit{one-hot}, donde se tiene solo una señal en alto por cada uno de los diferentes estados. Cada una de de estas señales se usó como señal de activación para su propio driver, cuyo esquemático se puede ver en la Figura \ref{fig:Driver Esquemático}. El resultado es una FSM que se mueve de manera secuencial entre cuatro estados cuándo esta recibe una señal de conteo síncrona con el reloj del sistema. Cada estado tiene su respectivo driver, el cuál permite accessar la alimentación de \qty{5}{V} solamente cuando su respectivo estado está activo.

Se definen los cuatro estados, junto a sus respectivas señales de driver, como $S0$, $S1$, $S2$, y $S3$. Las cuatro señales de salida del contador de la FSM, del bit menos significativo al bit más significativo, como $Q_0$, $Q_1$, $Q_2$, y $Q_3$. 

\begin{figure}[H]
    \centering
    \begin{center}\textit{(nueva foto del diagrama)}\end{center}
    \caption{Esquemático de drivers en FSM}
    \label{fig:Driver Esquemático}
\end{figure}


\begin{table}[H]
\centering
\caption{Desglose de codificación one-hot}
\label{tab:Drivers/Estados}
\begin{tabular}{cccccc}
\toprule
\textbf{Estado} & \textbf{Bits de Contador} & \textbf{Cuenta en FSM} & \textbf{Driver Activo} \\
\midrule
$S0$ & $\overline{Q_2}~\overline{Q_1}~ \overline{Q_0}$ & 0 & Driver 0 \\
$S1$ & $\overline{Q_2}~\overline{Q_1} ~Q_0$ & 1 & Driver 1 \\
$S2$ & $\overline{Q_2}~Q_1 ~\overline{Q_0}$ & 2 & Driver 2 \\
$S3$ & $\overline{Q_2}~Q_1 ~Q_0$ & 3 & Driver 3 \\
$S0$ & $Q_2~\overline{Q_1}~ \overline{Q_0}$ & 4 & Driver 0 \\
\bottomrule
\end{tabular}
\end{table}


El número preestablecido del contador se configuro para ser 0, de manera que todas sus salidas se niegan una vez que el pin de \textit{Preset Enable} recibe una señal en alto. En el sistema, \textit{Preset Enable} recibe la señal de salida del botón $\overline{Q}$, por lo que el contador de la FSM se encuentra perpetuamente en el estado $S0$ hasta que el botón sea activado. Una vez que el el botón se activa, su señal $\overline{Q}$ permanecerá baja, lu cuál le permitirá a la FSM cambiar a estados diferentes de $S0$.


La transición de un estado a otro es controlada por la señal $T_C$. Está sera explicada a fondo más adelante. Se recalca que el conteo del CD4029 se habilita mediante el pin de \textit{Carry-In}, el cuál está negado. Por lo tanto, la señal $T_C$ que recibe esta entrada usualmente se encuentra en alto, y se pone baja de manera síncrona durante un pulso del reloj cuendo se quiere que la FSM avance al siguiente estado. 

\subsection{Botón}

Como se mencionó anteriormente, el bóton cumple la función de controlar cuándo el contador de la FSM, así como los contadores en el subsistema de despliegue, se encuentran en su modo preestablecido. El diseño del botón está basado en el flip-flop SN7474, en conjunto con un buffer inversor SN7404. El uso de un flip-flop es crucial para poder guardar la "decisión" de apretar el botón hasta que la FSM regrese a su estado inicial. De esta manera, volver a apretar el botón durante un estado no deseado no tiene efecto alguno.

El 7474 se encuentra en modo asincrónico, por lo que no se utilizan sus pines de \textit{Clock} y \textit{Data}. Su pin de \textit{Set} está invertido. Conectarlo mediante una resistencia de \qty{10}{k\Omega} a $V_S$ permite que la entrada esté usualmente inactiva. Para activar el pin de \textit{Set}, se tienen dos botones conectados desde este pin a tierra. Estos proveen un switch momentario que pone la entrada del pin en bajo, que en torno fija las salida $Q=1$, y la salida $\overline{Q}=0$. Se evita un corto de $V_S$ a tierra mediante la resistencia de \qty{10}{k\Omega} mencionada.

El pin de \textit{Clear} también está negado, y se conecta a la señal del tercer bit de la FSM $\overline{Q_2}$, invertida mediante el 7404. Esto permite que cuando la cuenta de la FSM llegue a 4, se active el pin de \textit{Clear} del botón, y se fijen las salidas $Q=0$ y $\overline{Q}=1$. Cómo $\overline{Q}$ vuelve a estar en alto, esto activa los pines de \textit{Preset} de los contadores de la FSM y el subsistema de despliegue, preestableciendo las cuentas en 0 y 25 respectivamente. Cambiar el número del contador de la FSM de 4 a 0 no tiene un efecto notable en el sistema, ya que en ambos casos la máquina se encuentra en el estado $S0$, visto en la Tabla~\ref{tab:Drivers/Estados}.
\\
\\

Dado que, al encender el sistema por primera vez, el 7474 tiene la posibilidad de comenzar en cualquiera de sus dos estados, se tuvo que divisar una manera de que el pin de \textit{Clear} siempre reciba una señal baja en este momento. Para esto se utilizó otro de los inversores del 7404. Se conectó la entrada de un inversor a $V_S$ mediante un capacitor de \qty{10}{\mu F}, así como a tierra mediante una resistencia de \qty{10}{k\Omega}. La salida de este inversor se conectó al pin de \textit{Clear}.

Cuándo el circuito se enciende, este capacitor se encuentra descargado, por lo que actúa brevemente como un corto entre $V_S$ y la entrada del inversor. Esto permite que \textit{Clear} reciba una señal baja cuando el switch del subsistema de alimentación es cerrado. Posteriormente a este evento, el capacitor se logra cargar a los \qty{5}{V} de $V_S$, actuando como un abierto hasta que el sistema sea apagado. Una vez que este capacitor se carga, la entrada del inversor es fijada en bajo mediante la resistencia de \qty{10}{k\Omega}, lo cuál le envía perpetuamente una señal en alto al pin de \textit{Clear} y permite el funcionamiento ordinario del sistema. 





\subsection{Señal $T_C$}


\section{Estados} %Describir a qué está conectado el driver de cada estado (LEDs, cátodo, señal de audio, etc)

En esta sección se quiere mencionar cómo las señales tomadas de los 4 drivers interactúan con el resto del sistema. De esta manera, se puede tener el panorama completo de cómo funciona el sistema en su totalidad durante cada estado. 

\subsection{S0}



\subsection{S1}

\subsection{S2}

\subsection{S3}




%%%%%%%%%%%%%%%%%%%%%%%%%%%%%%%%%%%%%%%%%%%%%%%%%%%%%%%%%%%%%%
\begin{thebibliography}{99}
\bibitem{74ls74datasheet}
Texas Instruments, ``SN74LS74A Dual D-Type Positive-Edge-Triggered Flip-Flops With Preset and Clear Datasheet,'' Dallas, TX, USA. [Online]. Available: \url{https://c1555f5ec9.clvaw-cdnwnd.com/34662fcf1f1e607c561442431023ac8e/200000751-dcad7dda84/74LS74%20Datasheet.pdf}

\bibitem{7805datasheet}
onsemi, MC7800, MC7800A, MC7800AE, NCV7800 Positive 1.0 A Voltage Regulators Datasheet, Rev. 28, Apr. 2014. [Online]. Available: https://www.onsemi.com/pdf/datasheet/mc7800-d.pdf



\bibitem{electronicafp}
Electrónica FP, ``555 Astable,'' YouTube, 2021. [Online]. Available: \url{https://www.youtube.com/watch?v=OKDaA-yE2QY}
\bibitem{2n2222ds}
Philips Semiconductors, ``2N2222 NPN Switching Transistors,'' datasheet, Philips Semiconductors, Eindhoven, Netherlands. [Online]. Available: \url{https://www.alldatasheet.com/html-pdf/15067/PHILIPS/2N2222/995/4/2N2222.html}. [Accessed: May 2026].

\bibitem{tip31ds}
ON Semiconductor, ``TIP31, TIP31A, TIP31B, TIP31C: NPN Epitaxial Silicon Transistor,'' datasheet, ON Semiconductor, Phoenix, AZ, USA. [Online]. Available: \url{https://www.alldatasheet.com/html-pdf/499213/ONSEMI/TIP31/214/1/TIP31.html}. [Accessed: May 2026].

\bibitem{ne555art}
Ariat Tech, ``What is the NE555?'' \textit{Ariat Tech Blog}, 2024. [Online]. Available: \url{https://www.ariat-tech.es/blog/what-is-the-ne555.html#t2}. [Accessed: May 2026].


\bibitem{harris_digital}
D. M. Harris and S. L. Harris, \textit{Digital Design and Computer
Architecture}, 2nd ed. Morgan Kaufmann, 2013.


\end{thebibliography}

\end{document}
